\setchapterpreamble[u]{\margintoc}
\chapter{Random Variables and Distributions}
\labch{random-variables}

\sources{Bertsekas and Tsitsiklis, \emph{Introduction to Probability}, 2nd ed.\
\cite{bertsekas2008probability}, Chapters~2--3; MIT 6.041 \cite{ocw6041},
Lectures~5--12.}

A random variable attaches a number to each outcome, which turns questions about
events into questions about functions --- and therefore into questions that
calculus and linear algebra can answer.

\section{Discrete random variables}
\labsec{discrete-rv}

\begin{definition}
A \emph{random variable} is a function \(X:\Omega\to\R\). It is \emph{discrete}
when its range is finite or countable, and it is then described by its
\emph{probability mass function} \(p_X(x)=\Pr(X=x)\), which is non-negative and
sums to one.
\end{definition}

Three that recur:

\begin{description}
	\item[Bernoulli\((p\))] \(X\in\{0,1\}\) with \(\Pr(X=1)=p\). The model for a
	single binary label.
	\item[Binomial\((n,p\))] the number of successes in \(n\) independent
	Bernoulli trials, \(p_X(k)=\binom{n}{k}p^{k}(1-p)^{n-k}\).
	\item[Poisson\((\lambda\))] \(p_X(k)=e^{-\lambda}\lambda^{k}/k!\), the limit
	of Binomial\((n,\lambda/n)\) as \(n\to\infty\); counts of rare events.
\end{description}

\section{Continuous random variables}
\labsec{continuous-rv}

\begin{definition}
\(X\) is \emph{continuous} with \emph{density} \(f_X\ge0\) when
\[
	\Pr(a\le X\le b)=\int_a^b f_X(x)\,dx,
	\qquad
	\int_{-\infty}^{\infty}f_X(x)\,dx=1 .
\]
\end{definition}

A density is not a probability: \(f_X(x)\) may exceed one, and
\(\Pr(X=x)=0\) for every single point. What has meaning is \(f_X(x)\,dx\), the
probability of a small interval.

The \emph{cumulative distribution function} \(F_X(x)=\Pr(X\le x)\) works for both
kinds, is non-decreasing, and satisfies \(F_X'=f_X\) in the continuous case.

\begin{description}
	\item[Uniform\((a,b\))] constant density \(1/(b-a)\) on \([a,b]\).
	\item[Exponential\((\lambda\))] \(f_X(x)=\lambda e^{-\lambda x}\) for
	\(x\ge0\); memoryless waiting time.
	\item[Gaussian\((\mu,\sigma^{2}\))]
	\(f_X(x)=\dfrac{1}{\sqrt{2\pi\sigma^{2}}}\exp\Bigl(-\dfrac{(x-\mu)^{2}}{2\sigma^{2}}\Bigr)\).
\end{description}

\marginnote{The Gaussian earns its place twice over: \refch{limit-theorems}
explains why sums of many small effects look Gaussian, and \refch{estimation}
shows that assuming Gaussian noise turns maximum likelihood into least squares.}

\section{Joint, marginal, and conditional distributions}
\labsec{joint}

For two random variables the joint density \(f_{X,Y}\) determines everything.
Integrating out one variable gives a \emph{marginal},
\[
	f_X(x)=\int f_{X,Y}(x,y)\,dy,
\]
and dividing gives a \emph{conditional},
\(f_{Y\mid X}(y\mid x)=f_{X,Y}(x,y)/f_X(x)\) wherever \(f_X(x)>0\). Independence
is the statement \(f_{X,Y}=f_Xf_Y\).

\begin{remark}
Supervised learning is the business of modelling \(f_{Y\mid X}\) and nothing
else. The marginal \(f_X\) governs which inputs are seen, and a change in it
without a change in \(f_{Y\mid X}\) is the mildest form of the distribution shift
named in \refch{map-of-learning-problems}.
\end{remark}

\section{Transforming a random variable}
\labsec{transformations}

If \(Y=g(X)\) with \(g\) differentiable and invertible, the density transforms by
the Jacobian factor of \refsec{jacobian-determinant}:
\[
	f_Y(\vect{y})=f_X\bigl(g^{-1}(\vect{y})\bigr)\,
	\bigl|\det J_{g^{-1}}(\vect{y})\bigr| .
\]

\begin{example}
For \(X\sim\)~Gaussian\((0,1)\) and \(Y=\sigma X+\mu\), the inverse is
\((y-\mu)/\sigma\) with derivative \(1/\sigma\), giving the
Gaussian\((\mu,\sigma^{2})\) density. Standardising a feature is exactly this
change of variables run backwards.
\end{example}

\marginnote{Densities need the Jacobian; expectations do not. Computing
\(\E[g(X)]\) requires no such correction, as \refsec{lotus} explains --- a
distinction worth keeping firmly separate.}
